\section{Variance, standard deviation, and moments}

Expectation describes a center of balance of a distribution. It does not describe
how spread out the distribution is. Two continuous random variables can have the
same expectation and very different variability. To measure spread, we use
variance and standard deviation.

\subsection*{Deviation from the mean}

Let
\[
\mu=\E[X].
\]
The quantity
\[
X-\mu
\]
measures how far the random variable is from its mean. But the average deviation
\(\E[X-\mu]\) is always zero whenever the expectation exists, because
\[
\E[X-\mu]=\E[X]-\mu=0.
\]
Positive and negative deviations cancel each other. To measure size of deviation
without cancellation, we square it.

\begin{definitionbox}
Let \(X\) be a continuous random variable with density \(f_X\) and finite mean
\(\mu=\E[X]\). The \textbf{variance} of \(X\) is
\[
\Var(X)=\E[(X-
\mu)^2]=\int_{-\infty}^{\infty}(x-\mu)^2f_X(x)\,dx,
\]
provided the integral is finite. The \textbf{standard deviation} is
\[
\sigma_X=\sqrt{\Var(X)}.
\]
\end{definitionbox}

Variance is measured in squared units of \(X\). Standard deviation returns to the
original units, which often makes it easier to interpret.

\subsection*{Computational formula}

Variance can also be computed using
\[
\Var(X)=\E[X^2]-(\E[X])^2,
\]
provided the relevant expectations exist.

Indeed,
\[
\E[(X-
\mu)^2]=\E[X^2-2\mu X+
\mu^2].
\]
By linearity of expectation,
\[
\E[(X-
\mu)^2]=\E[X^2]-2\mu\E[X]+
\mu^2.
\]
Since \(\E[X]=\mu\), this becomes
\[
\Var(X)=\E[X^2]-2\mu^2+
\mu^2=\E[X^2]-
\mu^2.
\]

\begin{remarkbox}
The formula \(\Var(X)=\E[X^2]-(\E[X])^2\) is useful for computation, but the
definition \(\Var(X)=\E[(X-
\mu)^2]\) explains the meaning: variance is the
average squared distance from the mean.
\end{remarkbox}

\subsection*{Example: uniform distribution on \([0,1]\)}

Let \(X\) be uniformly distributed on \([0,1]\). We already know that
\[
\E[X]=\frac12.
\]
Also,
\[
\E[X^2]=\int_0^1 x^2\,dx=\frac13.
\]
Therefore
\[
\Var(X)=\frac13-\left(\frac12\right)^2=
\frac13-\frac14=\frac1{12}.
\]
The standard deviation is
\[
\sigma_X=\sqrt{\frac1{12}}=\frac1{2\sqrt3}.
\]

\subsection*{Example: uniform distribution on \([a,b]\)}

For \(X\sim \operatorname{Unif}(a,b)\),
\[
\E[X]=\frac{a+b}{2}.
\]
The second moment is
\[
\E[X^2]=\int_a^b x^2\frac1{b-a}\,dx
=\frac1{b-a}\left[\frac{x^3}{3}\right]_a^b
=\frac{b^3-a^3}{3(b-a)}.
\]
Using
\[
b^3-a^3=(b-a)(a^2+ab+b^2),
\]
we obtain
\[
\E[X^2]=\frac{a^2+ab+b^2}{3}.
\]
Hence
\[
\Var(X)=\frac{a^2+ab+b^2}{3}-\left(\frac{a+b}{2}\right)^2.
\]
After simplifying,
\[
\Var(X)=\frac{(b-a)^2}{12}.
\]
This formula shows that spread depends only on the length of the interval, not
on its location.

\subsection*{Example: triangular density}

Let
\[
f_X(x)=
\begin{cases}
2x, & 0\leq x\leq1,\\
0, & \text{otherwise}.
\end{cases}
\]
We previously found
\[
\E[X]=\frac23.
\]
Now
\[
\E[X^2]=\int_0^1 x^2\cdot2x\,dx=\int_0^1 2x^3\,dx=\frac12.
\]
Thus
\[
\Var(X)=\frac12-\left(\frac23\right)^2
=\frac12-\frac49=\frac1{18}.
\]
The density is concentrated toward the right side of the interval, so the mean is
larger than \(1/2\), while the variance reflects the remaining spread around
that shifted center.

\subsection*{Moments}

Moments summarize different aspects of a distribution. The \(k\)-th raw moment,
when it exists, is
\[
\E[X^k]=\int_{-\infty}^{\infty}x^k f_X(x)\,dx.
\]
The first raw moment is the expectation. The second raw moment helps compute
variance. Higher moments describe additional features of the shape of the
distribution.

Central moments are moments around the mean:
\[
\E[(X-
\mu)^k].
\]
The second central moment is variance. Higher central moments are used to study
asymmetry and tail behaviour, although this course mainly uses the first two
moments.

\subsection*{What variance does and does not say}

Variance measures average squared distance from the mean. A small variance means
that values tend to stay close to the mean. A large variance means that values
are more spread out. However, variance alone does not determine the entire shape
of a distribution. Different distributions may have the same mean and variance
but very different densities.

\begin{summarybox}
For a continuous random variable with density \(f_X\),
\[
\Var(X)=\int_{-\infty}^{\infty}(x-
\mu)^2f_X(x)\,dx,
\qquad
\mu=\E[X].
\]
Equivalently,
\[
\Var(X)=\E[X^2]-(\E[X])^2.
\]
Variance measures average squared spread around the mean; standard deviation is
its square root and is expressed in the original units.
\end{summarybox}
